\chapter{Conclusion}
\label{ch:conclusion}

This thesis started from a mismatch. Honeyword-based breach detection
depends on flatness --- the promise that an attacker who has read the
password file cannot identify the real credential in a sweetword set
with probability better than $1/k$ --- and every widely-deployed
honeyword generator has been evaluated against attackers that no
longer represent the state of the art. \textcite{dani2024passfilter}
showed that a small character-level CNN, trained as a binary
classifier on real-vs-decoy pairs, breaks flatness on essentially
every honeyword scheme they tested. The question this left open was
whether flatness can be restored against such an attacker without
redesigning the generator.

The answer we develop is affirmative but qualified. The same CNN
discriminator that PassFilter used as an attack,
reused unmodified as a defensive oracle at storage time, gates
sweetword sets against a small set of measurable flatness conditions.
Sets that fail are regenerated. Sets that pass are stored. The
generator itself is not modified, and neither is the attacker; only
the pipeline in which they are combined changes. On the ten-method
benchmark of \Cref{ch:results}, the two-stage filter with a
multiplicative pre-screening band reduces the attacker's $\mathrm{ASR}@5$
from around $0.75$ for the unfiltered baselines to $0.29$--$0.34$ for
the strongest configurations at $N_{\text{train}} = 10000$, with a
real password ranked near the middle of the twenty candidates
($F_{\text{rank}} \approx 0.42$--$0.46$). Under the thresholds adopted
in \Cref{sec:method-threshold-tradeoff}, these ratio filters accept
between $29\%$ and $78\%$ of passwords across corpus sizes --- well
inside the deployable range, and far above the single-digit acceptance
of the initial strict configuration.

\section{Summary of Contributions}

The concrete contributions of the thesis map to \Cref{ch:method}
and \Cref{ch:results}. First, an attacker-relative definition of
flatness formalized at the level of the attacker-induced distribution
over a sweetword set, characterized by four measurable quantities
($F_{\text{rank}}$, $F_{\text{ent}}$, top gap, and an
$\varepsilon$-flatness bound $q_{\max} \le 1/k + \varepsilon_{\max}$ on
the largest entry of the induced distribution). Second, a
hybrid flatness-filtering algorithm that combines candidate-level
pre-screening with set-level rejection sampling, generator-agnostic
and requiring no retraining of either the underlying generator or the
attacker.

The benchmark itself is the third contribution: ten honeyword
generation methods spanning rule-based tweaking,
representation-learning HoneyGen variants, and a self-trained GAN,
evaluated under a common experimental protocol at three training sizes,
and against three attacker classes --- the
HoneyGen-tuned CNN, a mixed-corpus CNN trained on several decoy
families, and a frequency ranker --- which together separate genuine
flatness from single-attacker artifacts.

\section{What the Numbers Actually Say}

Three quantitative findings deserve to be extracted from
\Cref{ch:results} and stated cleanly enough for a reader in a hurry.

Unfiltered honeyword generators are highly distinguishable to the
CNN. All four baselines (Tweaking, Password-Model, Hybrid-Model,
HoneyGen Baseline) report $\mathrm{ASR}@5$ in the range $0.68$--$0.77$
across every training-set size we tested, about three times the
random-guessing rate $0.25$.

Flatness filtering reduces $\mathrm{ASR}@5$ substantially, and the
ratio band is what does most of the work. At $N_{\text{train}} = 10000$
the Prescreen-Ratio and Hybrid-Ratio variants push $\mathrm{ASR}@5$
down to $0.341$ and $0.293$ with $F_{\text{rank}}$ near the ideal
middle value of $0.5$ --- genuine near-uniform ranking --- while the
additive filters achieve moderate reductions to the $0.51$--$0.61$
range. The relative ordering of methods is preserved as the training
corpus varies over $2000$, $5000$, and $10000$. Because acceptance is
treated as a first-class deployability requirement
(\Cref{sec:results-cost}), the recommended configurations are the
ratio filters read through a two-stage criterion: the flattest methods
among those that remain deployable.

The GAN's apparent flatness is a distributional artifact and should not
be read as a defence. Following the CNN's top guesses, its
$\mathrm{ASR}@5$ at $N_{\text{train}} = 10000$ is essentially zero, but
$F_{\text{rank}}$ is $0.920$: the real password is ranked near position
twenty on the average set. The CNN considers the GAN-generated
candidates to look \emph{more} like real passwords than the actual real
password does, so a rational attacker inverts the ranking and reaches
$\mathrm{ASR}@5 = 0.872$ --- the value the two-sided rational-attacker
floor reports, which makes the GAN one of the weakest methods in the
benchmark rather than the strongest. Two independent checks confirm the
bias rather than flatness: under the frequency attacker the same GAN
sets report $\mathrm{ASR}@5 = 0.388$, the highest in that column, and
under the mixed-corpus attacker the GAN's $\mathrm{ASR}@5$ rises from
about $0.07$ to $0.99$ at $N_{\text{train}} = 5000$. This artifact is
pronounced at the larger corpus sizes; at $N_{\text{train}} = 2000$ the
GAN's $F_{\text{rank}}$ is closer to the middle, so the inversion is a
property that sharpens as the generator is trained on more data.

\section{Limitations}
\label{sec:conclusion-limitations}

The framework has five limitations that a deployment would need to
reason about explicitly.

\paragraph{Coverage below 100\%.}
The most important limitation for real-world use is that the
countermeasure does not protect every user. The filter stores a
sweetword set only when it can certify that set as flat; when it cannot
do so within the regeneration budget, it gives up on that password, and
that account receives no flatness-guaranteed set. At the operating point
of this thesis the strongest (ratio) filters certify only between
$29\%$ and $78\%$ of passwords (\Cref{tab:accept}), which leaves a large
share of users --- in the worst case the majority --- without the
guarantee. A deployable system, however, has to protect $100\%$ of
accounts: it cannot simply refuse to store honeywords for the users it
finds difficult. As it stands the method would therefore need to be
paired with a fallback for the rejected passwords (for instance an
unfiltered set), which weakens the guarantee for exactly those users.
Reaching near-total coverage without sacrificing flatness is the single
biggest obstacle to real-world adoption, and we return to it as an open
question below.

\paragraph{Model-relative guarantees.}
The filter enforces flatness against one specific attacker: the
character-level CNN of \Cref{sec:method-cnn}. A different attacker ---
a larger CNN, a transformer, a discriminator that uses side
information beyond the password string --- may recover
distinguishability on sets the filter has certified. The guarantees
are best read as flatness against this class of learned attacker,
not as flatness in general. \Cref{sec:results-freq} shows that the
guarantees transfer partially to the frequency attacker, which is
weaker; whether they transfer to stronger models is an empirical
question that requires running each candidate attacker through the
same benchmark.

\paragraph{One-sided set-level acceptance.}
The set-level acceptance test constrains only the maximum of the
induced distribution $q$. A set in which some honeyword receives an
anomalously low $q$ would pass the test even though $q$ is not close
to uniform. In our data this does not translate into ASR degradation,
because ASR responds to the top of the ranking rather than the tail,
but the acceptance predicate could be tightened to bound both ends of
$q$. This is a small extension of \Cref{alg:hff} and would be worth
doing in a production build of the filter.

\paragraph{Acceptance-rate cost of the ratio filter.}
The strongest security in the benchmark, from the Prescreen-Ratio and
Hybrid-Ratio filters, comes at the lowest acceptance rates among the
filtered methods --- between $29\%$ and $78\%$ across corpus sizes,
against $65\%$--$97\%$ for the additive filters --- and a higher
average number of regeneration attempts, up to about $4.6$ per
accepted set (\Cref{tab:accept}). Under the relaxed thresholds this is
no longer the crippling single-digit acceptance of the initial
configuration: every filtered method clears the $25\%$ deployability
mark at every training size. The residual cost is nonetheless real and
falls entirely at storage time, since filtering happens once per
account for the lifetime of the credential. A deployment that cannot
afford the ratio filter's rejection budget can move to an additive
filter and recover most of the security gain over the unfiltered
baseline at a substantially higher acceptance rate.

\paragraph{Corpus and language.}
All experiments use the cleaned RockYou corpus, a leak of primarily
English-language, user-chosen web passwords from a specific consumer
service. The distributional properties of this corpus --- long
low-probability tails, common suffixes, digit sequences at the end
--- are what the FastText embeddings, the GAN, and the CNN all learn
to model. Whether the same numbers hold on corpora with different
structural properties (enterprise credentials, non-English leaks,
machine-generated passwords) is an empirical question. Extending the
benchmark to \emph{darkc0de} as a second corpus is on the immediate
next-steps list for this pipeline.

\section{Open Questions}

Beyond the immediate limitations, five research questions remain open.

The first is adversarial adaptation. The filter, as constructed, is a
static defense: the oracle CNN is trained once and reused. An attacker
who knew the filter existed and had access to a sample of accepted
sets could train against those sets rather than against the raw
generator output. Whether the generator, the oracle, and the attacker
can be co-trained toward a stable equilibrium at
$\mathrm{ASR}@5 \approx 0.25$ is a question this thesis does not
resolve. The pieces are in place to run the experiment; the answer
requires it.

The second is transferability across a class of attackers. The
current guarantees are against one CNN. A more useful guarantee would
be against a class of models --- for example, any CNN of bounded
depth and width, or any linear classifier over a fixed feature set.
The filter could then use an ensemble oracle that discriminates
against the worst-case member of the class.
\textcite{chakraborty2022hbat} discuss related notions in the survey,
and connecting them to the flatness formalism of
\Cref{sec:method-formalism} is a natural extension.

The third is the theoretical relationship between the flatness
metrics. $F_{\text{rank}}$, $F_{\text{ent}}$, and $\varepsilon$-flatness
are correlated but not equivalent, and the acceptance predicate uses
all three plus the top gap. It is not obvious that this is the
smallest set of conditions sufficient to bound $\mathrm{ASR}@t$, or
that all three are individually necessary. A cleaner theoretical
characterization --- an inequality connecting $q_{\text{real}}$,
$F_{\text{ent}}$, and expected $\mathrm{ASR}@t$ --- would allow the
filter to constrain fewer quantities without giving up guarantees.

The fourth is the LLM-based generator line. The construction of
\Cref{sec:method-filter} is generator-agnostic and would apply
without modification to sweetword sets produced by the
chunk-manipulation and Metropolis-Hastings large language model of
the kind reviewed in \Cref{sec:related-llm}. Whether the resulting
$\mathrm{ASR}@5$ curves match, exceed, or fall below the
ratio-filtered HoneyGen numbers reported here is an experiment worth
running.

The fifth is coverage. As the limitations above stress, the filter
protects only the passwords for which it can certify a flat set, and at
the current operating point that fraction is well below $100\%$. A
practical countermeasure must cover every account, so the open question
is how to reach near-total coverage without weakening flatness. Possible
routes include stronger or password-specific generators that produce
flat candidates more reliably, thresholds adapted per password to the
hardest cases, or a principled fallback for the residual passwords that
no generator can flatten. Which of these --- or which combination ---
closes the gap is the most important question this thesis leaves open.

\section{Closing}

Honeywords are a modest defense with an outsized property: they turn
a silent failure --- offline password recovery from a stolen hash
file --- into an observable event at the authentication server. That
property is worth preserving as attacker capabilities improve, and
the argument this thesis makes is that each new class of attacker
should be reused as an enforcement tool at storage time. The
classifier that would otherwise recover the real password becomes
the gatekeeper that prevents any sweetword set with a recoverable
real password from being stored. The reversal is small and the
numbers are moderate, but the direction feels right: as long as
attackers keep getting better, defenders can keep the same tools,
pointed the other way.
